\pageheading{EXACT LABORATORY\; /\; COEFFICIENTS}{Integer coordinates for the nonlinear family}
Normalize the Taylor coefficients by
\[
A_n(q,r)=(2n+1)!c_n(q,r),\qquad A_0=1.
\]
The recurrence becomes
\[
\boxed{A_n=\sum_{i=0}^{n-1}\binom{2n}{2i+1}
q^{2i+1}r^{2(n-1-i)+1}A_iA_{n-1-i},\quad n\ge1.}
\]
Therefore $A_n\in\mathbb Z[q,r]$ has nonnegative coefficients and is
symmetric in its two variables. The first coordinates are
\[
\begin{aligned}
A_1&=2qr,\\
A_2&=8q^2r^2(q^2+r^2),\\
A_3&=q^3r^3[48(q^2+r^2)(q^4+r^4)+80q^2r^2].
\end{aligned}
\]
At $(q,r)=(1,1)$ these are the tangent numbers
$1,2,16,272,7936,\ldots$, the known OEIS A000182 anchor.
At $(1,2)$ the exact sequence begins
$1,4,160,35200,38041600,\ldots$.
\smallfig{17_tree_addresses}{Ordered trees retain left/right path information.
The two displayed shapes contribute different monomials to the same
factorial-normalized coefficient. Their labels are exact tree data.}
At rational scales, an integer denominator clearing is available. If $D$
clears every contraction denominator and $H$ clears every bank weight,
\[
H^nD^{n(n+1)}A_n\in\mathbb Z.
\]
The maximum root-path exponent sum over a full binary tree is $n(n+1)$.
The clearing provides exact sequence exports without implying a new OEIS entry.

\pageheading{EXACT LABORATORY\; /\; REGIMES}{A formal sequence is not always an analytic germ}
\smallfig[0.68]{18_parameter_regimes}{The full nonnegative real parameter map separates
finite-order entire functions, infinite-order entire edges, tangent, and
zero-radius formal series. Zero scales give the linear function.}
If $q>0$, $r>1$, the recurrence includes its $i=0$ term and gives
\[
c_n\ge\frac{q^nr^{n^2}}{3\cdot5\cdots(2n+1)}.
\]
The root test proves zero radius. By symmetry this covers either scale above
one, provided the other is positive. The exported positive integer asymmetric
rows belong to this formal regime.

If $0<q<1$, $r=1$, the equation is linear in the uncontracted value:
$T'=1+T(qz)T(z)$. A known coefficient on $|z|<R/q$ gives a linear ODE
there, so disk bootstrapping proves entire extension. On the positive ray,
$(\log T)'=T(qR)+1/T(R)\ge qR$, giving an exponential lower bound for
$T$. Substitution back gives, for $R\ge2/q$,
\[
\log T(R)\ge\frac R2\exp(q^3R^2/8-q/2).
\]
Thus the entire order is infinite. The corner $(1,1)$ is tangent with poles.
These exact boundaries protect the analytic interpretation of the catalog.

\pageheading{EXACT LABORATORY\; /\; PROJECTORS}{Polyphase sections preserve the full sequence}
For an ordinary generating function $G(z)=\sum_{m\ge0}a_mz^m$ and
integer $b\ge1$, write
\[
G(z)=\sum_{j=0}^{b-1}z^jG_j(z^b),\qquad
G_j(t)=\sum_{n\ge0}a_{bn+j}t^n.
\]
The root projector is
\[
P_{b,j}G(z)=\frac1b\sum_{k=0}^{b-1}\omega^{-jk}G(\omega^kz)
=z^jG_j(z^b),\qquad\omega=e^{2\pi i/b}.
\]
The same projector works on an EGF, but the section retains its original
factorials: $\sum_na_{bn+j}z^{bn+j}/(bn+j)!$.
Removing those denominators changes the representation and requires an
explicit transport. Root filtering does not do it automatically.

For causal FIR taps $h_\ell$, $b_m=\sum_\ell h_\ell a_{m-\ell}$ with
zero extension to negative indices, split $A(z)$ and $B(z)$ into $b$ phases.
Then $B_i(t)=\sum_jH_{ij}(t)A_j(t)$, where
\[
\boxed{H_{ij}(t)=\sum_{\substack{\ell\ge0\\j+\ell\equiv i\pmod b}}
h_\ell t^{(j+\ell-i)/b}.}
\]
The exponent is the nonnegative carry to the next phase block. This is an
exact coefficient law, with no approximation or loss of terms.

\pageheading{EXACT LABORATORY\; /\; TAPS}{The carry matrix is part of the filter}
\smallfig{16_polyphase}{A three-phase filter with taps $[1,2,-1]$.
The matrix carries terms across block boundaries. Interleaving its output
reproduces the direct scalar FIR calculation exactly.}
With these taps and $b=3$,
\[
H(t)=\begin{pmatrix}
1&-t&2t\\2&1&-t\\-1&2&1
\end{pmatrix}.
\]
The entries arise from the same congruence condition. For instance an input
in phase two shifted by one lands in phase zero of the next block, contributing
$2t$ to $H_{0,2}$.

\claim{Three operations that must be kept distinct}
For an EGF $E_a(z)=\sum_ma_mz^m/m!$, differentiation transports future
coefficients: $E_a^{(k)}=\sum_ma_{m+k}z^m/m!$.
Causal ordinary taps instead act by repeated zero-constant integration:
\[
E_b=\sum_\ell h_\ell I^\ell E_a.
\]
Multiplication of two EGFs gives Hurwitz convolution,
$b_m=\sum_{k=0}^m\binom mk a_kh_{m-k}$.
It is not ordinary sequence FIR convolution. Intrinsic nonlinear feedback
taps alter the defining family and are a fourth operation with their own
coefficient recurrence.

\pageheading{EXACT LABORATORY\; /\; REUSABLE CODE}{A catalog with explicit provenance}
The standard-library \texttt{family\_lab.py} exposes exact coefficient
polynomials, normalized sequences, positive-bank coefficients, rational
denominator clearing, root sections, Hurwitz products, ordinary FIR taps,
and polyphase filter matrices. The catalog contains 18 rows through index 12.

Every row records parameters, coefficient normalization, arithmetic regime,
and whether its entries are integers or rationals. Eight positive integer
asymmetric specializations are labeled zero-radius formal sequences. Entire
interior examples and an infinite-order edge are labeled separately.

\begin{tabular}{@{}p{0.27\linewidth}p{0.35\linewidth}p{0.30\linewidth}@{}}
\toprule
Example & First normalized coordinates & Interpretation\\\midrule
$(1,1)$ & $1,2,16,272,7936$ & Known tangent anchor\\
$(1,2)$ & $1,4,160,35200,38041600$ & Integer, formal radius zero\\
$(1/2,1/2)$ & $1,1/2,1/4,1/8,1/16$ & Scaled sinh, entire order one\\
$(1/3,2/3)$ & $1,4/9,160/729,18560/177147,\ldots$ & Asymmetric entire order one\\
\bottomrule
\end{tabular}

\claim{Implementation boundaries are enforced}
Phase residues and bases must be integers, and a phase bank cannot be empty.
The independent review found that fractional residues and zero phases were
previously accepted. Both domain-validation defects were fixed and covered
by regression assertions before release.

The 8,038 exact API checks compare separate constructions: Picard integration,
sine/cosine tangent division, root filters in a cyclotomic quotient,
direct and phase-matrix taps, and coefficient-level antiderivatives.
They verify finite implementation ranges. The algebraic laws above have
separate proofs in Appendix E.

The catalog is ready for comparison with OEIS, with known anchors identified.
It makes no claim that an unmatched prefix establishes a new sequence, and
no sequence was submitted externally as part of this release.
